\section{Bulk Loader}
\label{arch:bulkloading}

Bulk loading inserts a large number of records into the database from an external source.
Systems typically disable constraint checks, index updates, and logging during bulk loading to maximize throughput.
Graph databases must support bulk loading, because graphs can be large (the largest Graphalytics graphs have billions of edges) and transactional insertion at such scale is prohibitively slow.

\subsection{Preprocessing Pipeline}
\label{arch:bulk-preprocess}

Before loading, the raw edge list must be transformed into a format compatible with \project{}'s representations.
The input is a CSV edge file sorted by source vertex ID.
A single C++ program replaces the original multi-stage shell pipeline (\texttt{awk}, \texttt{sort}, \texttt{split}) with a two-pass design using in-process parallel sorting.

\paragraph{Pass 1: Forward edges and reverse-edge chunking.}
The program reads the CSV edge file in a single sequential pass.
Forward edges arrive sorted by source ID, so the program groups them by source vertex and writes them directly to binary adjacency list partition files.
Simultaneously, it reverses each edge $(u, v)$ to $(v, u)$ and buffers it in memory.
When the buffer reaches a configurable size (default 10\,GB, approximately 625M edges), the chunk is sorted in parallel using \texttt{tbb::parallel\_sort} and flushed to a binary run file on disk.

\paragraph{Pass 2: K-way merge of reverse edges.}
A min-heap-based $k$-way merge combines the reverse-edge run files.
The program groups the merged stream by source vertex and writes it to binary adjacency list partition files for the in-neighborhood tables.

\paragraph{Output format.}
The output uses a binary adjacency list format:
\[
  \texttt{[node\_id}_{64} \;|\; \texttt{degree}_{32} \;|\; \texttt{dst\_id}_1 \;|\; \cdots \;|\; \texttt{dst\_id}_{d}\texttt{]}
\]
where each field is written in its native binary representation.
This format eliminates text parsing overhead and is approximately 50\% smaller than the equivalent ASCII representation.

\paragraph{Partitioning.}
Output files are partitioned by vertex ID range: $\text{partition}(v) = \lfloor v \cdot N / |V| \rfloor$, where $N$ is the number of partitions (typically matching the core count).
Because each adjacency list is assigned to exactly one partition by its source vertex, no adjacency list spans a partition boundary, and no cross-partition merge pass is required.

\subsection{Transactional Loader}
\label{arch:bulk-txn}

The transactional loader uses standard \wt{} cursor inserts within implicit transactions.
Each thread opens a separate session and cursor to the target tables, reads its assigned shard of the preprocessed edge list, and inserts records line by line.
Because no \texttt{bulk} configuration is used, multiple threads can insert into the same table concurrently.
The transactional loader is simpler to implement and useful for small-to-medium graphs, but it is slower than the bulk-cursor approaches below due to transaction overhead and write-ahead logging.

\subsection{Standard \wt{} Bulk Loader}
\label{arch:bulk-wt}

\wt{} provides an option to bulk load data into a \emph{newly created} table by opening a cursor with the \texttt{bulk} configuration parameter.
The bulk cursor bypasses transaction overhead and writes pages directly, but imposes two constraints: (1)~keys must be inserted in sorted order, and (2)~only one bulk cursor may be open on a table at a time, precluding multi-threaded insertion into the same table.

\project{} uses the standard bulk loader by streaming presorted partition files through a single bulk cursor per table.
Each table is loaded sequentially so that the full \wt{} cache is available during its pass, maximizing page-write throughput.
The bulk cursor accepts both text and binary adjacency list formats via a templated reader interface.
This approach serializes insertion into any individual table, but the preprocessing step (\autoref{arch:bulk-preprocess}) ensures that data arrives in the required sorted order.

\subsection{Modified \wt{} Bulk Loader}
\label{arch:bulk-modified}

To remove the single-threaded constraint, \project{} uses a modified \wt{} bulk loader that enables multiple threads to insert into the same table concurrently through bulk cursors.
The modification operates in two phases:

\paragraph{Phase 1: Parallel leaf-page creation.}
The modification removes the exclusive-access requirement (\code{WT\_DHANDLE\_EXCLUSIVE}) from the bulk cursor open path, allowing multiple threads to open bulk cursors on the same table simultaneously.
Each thread receives its own reconciliation context and a private backing file.
Threads write leaf pages independently to their backing files, recording the key and page-address for each leaf page written.
A spinlock-protected counter assigns file IDs atomically.
Because each thread writes only to its own backing file, the leaf-page creation phase requires no inter-thread synchronization beyond the file ID allocation.

\paragraph{Phase 2: Hierarchical internal-page merge.}
When a thread finishes inserting its partition, it reads its backing file and constructs internal B-tree pages one level at a time from the recorded key--address pairs.
Each thread then contributes its final root-level key--address pairs to a shared array, protected by a spinlock.
The last thread to complete (detected via an atomic decrement of the open-cursor count) sorts all root-level keys and builds the single root page of the B-tree.

This design maximizes parallelism during the I/O-intensive leaf-page writing phase, with synchronization limited to spinlock-protected metadata operations during the merge phase.
The modifications were originally built against \wt{} 3.2.1 and have since been ported to the \wt{} 11.0.0 release used by the rest of \project{}.

\subsection{Comparison}
\label{arch:bulk-comparison}

Table~\ref{tab:bulk-load-times} compares the three bulk loaders on the Twitter MPI graph (52.6\,M vertices, 1.96\,B directed edges, 16 partitions).
The test machine has a 14-core Intel Xeon W-2275 at 3.30\,GHz (28 hardware threads), 128\,GB of RAM, and a rotational hard disk.
The \wt{} cache is 10\,GB.
Each run starts from a cold page cache and loads the graph into both the adjacency-list and split edge-key representations.

\begin{table}[t]
\centering
\begin{tabular}{lrrr}
\toprule
Phase & Transactional & Std.\ Bulk & Par.\ Bulk \\
\midrule
Degree collection    & ---           & 261 & 247 \\
Adjacency lists      & ---           & 301 &  33 \\
Edge table           & ---           &  53 &  81 \\
Edge-key tables      & ---           & 159 & 150 \\
Nodes                & ---           &  18 &  15 \\
\midrule
\textbf{Total}       & $>$3{,}600    & \textbf{797} & \textbf{535} \\
\bottomrule
\end{tabular}
\caption{Bulk loading times (seconds) for the Twitter MPI graph using 16 threads.
The transactional loader did not finish within the 60-minute timeout.}
\label{tab:bulk-load-times}
\end{table}

The transactional loader did not complete within a 60-minute timeout.
Forward-edge insertion took 1{,}314 seconds across 16 threads, reverse-edge insertion took 974 seconds, and the node-insertion pass was still running when the timeout expired.
Each thread inserts records into both database representations through standard cursors within implicit transactions.
Transaction overhead, write-ahead logging, and the random-access pattern of node insertion account for the cost.

The standard \wt{} bulk loader completed in 797 seconds.
A degree-collection pass (261\,s) reads all 32 partition files to build in-degree and out-degree arrays before insertion begins.
Adjacency-list streaming through a single bulk cursor took 301 seconds for the in- and out-neighbourhood tables combined.
The edge and edge-key tables took 212 seconds, and node insertion took 18 seconds.

The modified bulk loader finished in 535 seconds, a $1.5\times$ speedup over the standard bulk loader.
Adjacency-list streaming dropped from 301 to 33 seconds ($9.1\times$) because 16 threads write leaf pages to independent backing files with no inter-thread synchronization.
The degree-collection pass (247\,s versus 261\,s) showed little change because it reads partition files sequentially regardless of thread count.
Edge-table streaming was slightly slower in the parallel case (81\,s versus 53\,s) because the hierarchical merge that combines each thread's B-tree into a single root adds work that offsets the parallel leaf-page writes for fine-grained per-edge keys.

All three loaders read the same preprocessed partition files, so the preprocessing time (\autoref{arch:bulk-preprocess}) is shared.
Table~\ref{tab:preprocess-times} breaks down the two preprocessing passes on the same machine with 16 threads.
Pass~1 is dominated by the sequential CSV read and forward-adjlist writes; the two \texttt{tbb::parallel\_sort} calls on the reverse-edge chunks account for only 16 of the 310 seconds.
Pass~2, a two-way merge of the sorted runs, completes in under a minute.
Including preprocessing, the end-to-end time from raw edge list to \wt{} database is 1{,}164 seconds for the standard bulk loader and 902 seconds for the modified bulk loader.

\begin{table}[t]
\centering
\begin{tabular}{lr}
\toprule
Step & Time (s) \\
\midrule
Pass 1: CSV read + forward adjlists + reverse sort & 310 \\
\quad Chunk 0 sort (1.34\,B edges)                 &  11 \\
\quad Chunk 1 sort (621\,M edges)                   &   5 \\
Pass 2: $k$-way merge of reverse runs              &  57 \\
\midrule
\textbf{Total preprocessing}                        & \textbf{367} \\
\bottomrule
\end{tabular}
\caption{Preprocessing times (seconds) for the Twitter MPI graph using 16 threads.
Pass~1 reads the CSV edge file, writes forward adjacency lists, and sorts reverse-edge chunks.
Pass~2 merges the sorted chunks into reverse-adjacency partitions.}
\label{tab:preprocess-times}
\end{table}